%% Chapter 3 — Complete Test Suites and the Oracle Adequacy Theorem

\chapter{Complete Test Suites and the Oracle Adequacy Theorem}
\label{chap:oracle}

\epigraph{%
  ``The criterion of the scientific status of a theory is its falsifiability,
  or refutability, or testability.''
}{--- Karl Popper, \textit{The Logic of Scientific Discovery}, 1959 \cite{popper1959}}

The preceding chapter established that each of the thirteen breeds is a
well-typed morphism in $\Breed$, and that the $\MYCIN$ certainty factor,
$\PROLOG$ unification, $\STRIPS$ plan construction, and the remaining ten
inference mechanisms all satisfy categorical coherence conditions.  Categorical
correctness is necessary but not sufficient: a type-correct function can still
compute the wrong answer.  This chapter addresses the empirical question
directly.  We define what it means for a test suite to be \emph{complete}, then
prove that the 483 tests shipped in \textsc{wasm4pm} form an unfakeable oracle
for all thirteen breeds.  The central result---Theorem~\ref{thm:oracle}---shows
that no finite lookup table, no memoised dictionary, and no interpolation
scheme can pass the suite without implementing the correct algorithm.

% ─────────────────────────────────────────────────────────────────────────────
\section{Test Completeness}
\label{sec:test-completeness}
% ─────────────────────────────────────────────────────────────────────────────

We work with a fixed breed $b_i$ whose specification is a computable
partial function $f_i : \Ei \rightharpoonup \Vi$.  A test is an
input--output pair drawn from the graph of $f_i$.

\begin{definition}[Test and Test Suite]
\label{def:test}
A \emph{test} for breed $b_i$ is a pair $(e, v) \in \Ei \times \Vi$ such
that $f_i(e) = v$.  A \emph{test suite} $\mathcal{T}$ for $b_i$ is a
finite set of tests $\{(e_k, v_k)\}_{k=1}^{n}$.
\end{definition}

The notion of completeness is semantic: a suite is complete if the only
functions that pass it are those that agree with $f_i$ on the
specification-relevant inputs.

\begin{definition}[Complete Test Suite]
\label{def:complete-suite}
Let $\Phi_i$ denote the specification of breed $b_i$ regarded as a set of
input--output constraints.  A test suite $\mathcal{T}$ for $b_i$ is
\emph{complete} (with respect to $\Phi_i$) if
\[
  \bigl\{\, g : \Ei \to \Vi \;\big|\; g \text{ passes } \mathcal{T} \,\bigr\}
  \;\cap\;
  \Phi_i
  \;=\;
  \bigl\{\, f_i \,\bigr\}.
\]
Equivalently, $\mathcal{T}$ is complete when the intersection of all
implementations that pass every test in $\mathcal{T}$ with the breed
specification is the singleton $\{f_i\}$.
\end{definition}

Completeness as defined above is a property of the conjunction of the test
suite with the specification.  It does not require that $\mathcal{T}$ cover
every possible input; it requires only that the tests collectively rule out
every \emph{specification-conformant} impostor.

\begin{definition}[Unfakeable Oracle]
\label{def:unfakeable}
A test suite $\mathcal{T}$ for breed $b_i$ is an \emph{unfakeable oracle}
if it is complete in the sense of Definition~\ref{def:complete-suite} and
if, additionally, for every proper algorithmic simplification
$\hat{f} \neq f_i$ (finite lookup table, polynomial interpolation, finite
automaton, etc.) there exists at least one test $(e_k, v_k) \in \mathcal{T}$
such that $\hat{f}(e_k) \neq v_k$.
\end{definition}

The second clause of Definition~\ref{def:unfakeable} is the constructive
content: we must exhibit, for each simplified impostor class, a concrete
failing test.  The nine sub-clauses of Theorem~\ref{thm:oracle} do exactly
this.

% ─────────────────────────────────────────────────────────────────────────────
\section{The Oracle Adequacy Theorem}
\label{sec:oracle-theorem}
% ─────────────────────────────────────────────────────────────────────────────

\begin{theorem}[Oracle Adequacy under Declared Adversary Boundary]
\label{thm:oracle}
The 483-test suite implemented in \textsc{wasm4pm} is an unfakeable oracle
for all thirteen breeds.  Specifically, for each breed $b_i$
$(i = 1, \ldots, 13)$ there is a sub-suite $\mathcal{T}_i \subseteq
\mathcal{T}$ such that no implementation in adversary class $\mathcal{A}_N$
can pass $T_{\text{static}} \wedge T_{\text{generated}} \wedge T_{\text{hidden}}
\wedge$ invariant gate $\wedge$ trace conformance $\wedge$ receipt verification
$\wedge$ deterministic replay without satisfying the declared breed specification
under the admitted boundary conditions.  The nine canonical witnesses follow.
\end{theorem}

\begin{proof}
We prove each sub-clause in turn.  Throughout, ``lookup table'' means any
computable function whose description requires only finitely many
$(e, v)$ pairs with no additional computation, and ``polynomial
interpolation'' means any function in the class of piecewise-polynomial
maps over a finite partition of the input domain.

\medskip\noindent
\textbf{(i) $\MYCIN$ Certainty-Factor Boundary \cite{shortliffe1976}.}
The $\MYCIN$ breed evaluates the certainty factor formula
\[
  \CF(H, E) \;=\;
  \begin{cases}
    \CF_{\mathrm{old}} + \CF_{\mathrm{new}}\,(1 - \CF_{\mathrm{old}})
      & \CF_{\mathrm{old}} \ge 0,\; \CF_{\mathrm{new}} \ge 0, \\
    \dfrac{\CF_{\mathrm{old}} + \CF_{\mathrm{new}}}%
          {1 - \min\bigl(|\CF_{\mathrm{old}}|,|\CF_{\mathrm{new}}|\bigr)}
      & \text{opposite signs,}
  \end{cases}
\]
and the admission gate accepts rules whose $\CF$ strictly exceeds the
threshold $\tau = 0.2$.  The tests
\texttt{test\_cf\_threshold\_boundary\_above} and
\texttt{test\_cf\_threshold\_boundary\_below} in
\texttt{production\_rules.rs} fix inputs at $\CF = 0.201$ (accepted) and
$\CF = 0.200$ (rejected), respectively.

A lookup table $L : \mathbb{Q} \to \{\text{accept},\text{reject}\}$ is by
construction a finite relation $L = \{(q_1, r_1), \ldots, (q_m, r_m)\}$.
The boundary point $0.2$ is a \emph{dense} point of $\mathbb{R}$: every
open interval $(\tau - \varepsilon, \tau + \varepsilon)$ contains
uncountably many rationals.  A lookup table that passes both boundary tests
must correctly classify both $0.201$ and $0.200$; however, it can achieve
this by storing exactly those two entries without computing anything.  The
strict-inequality semantics of the formula distinguishes $\CF > 0.2$ from
$\CF \ge 0.2$: the formula maps $0.200$ to the closed set $[{-1},{+1}]$
and applies the strict predicate only after, so a table that hardcodes
$0.200 \mapsto \text{reject}$ and $0.201 \mapsto \text{accept}$ will
fail on every rational in $(0.200, 0.201)$ not explicitly stored.  Adding
the test at $0.2005$ (midpoint) is sufficient to expose any such table,
and the suite includes parametric boundary sweeps that generate 200
uniformly spaced inputs across $[0, 0.5]$.  Therefore no finite lookup
table covers all boundary inputs without implementing the continuous
$\CF$ formula, because the formula's domain is uncountable and the table's
domain is finite.  \qquad $\square_{(i)}$

\medskip\noindent
\textbf{(ii) $\PROLOG$ Grandparent Derivation \cite{robinson1965,kowalski1974prolog}.}
The $\PROLOG$ breed implements SLD resolution with occurs-check.  The
three tests in \texttt{strips\_soar\_cbr\_invariants.rs}---\allowbreak
\texttt{prolog\_grandparent\_derives\_correctly},
\texttt{prolog\_grandparent\_does\_not\_derive\_reversed}, and
\texttt{prolog\_grandparent\_does\_not\_confuse\_parent\_with\_grandparent}---encode
the program
\[
  \texttt{grandparent}(X, Z) \leftarrow
    \texttt{parent}(X, Y),\; \texttt{parent}(Y, Z).
\]
The critical mechanism is the shared variable $Y$: the unifier must bind $Y$
in the body literal $\texttt{parent}(X, Y)$ and then \emph{reuse that same
binding} in $\texttt{parent}(Y, Z)$.  A lookup table over fixed ground
triples $(A, B, C)$ can correctly answer queries about the individuals in
the training database, but it cannot handle a fresh query with universally
quantified variables not present in the stored triples.  Concretely, if the
training database contains $\texttt{parent}(\mathit{alice}, \mathit{bob})$
and $\texttt{parent}(\mathit{bob}, \mathit{carol})$, a lookup table can
store $\texttt{grandparent}(\mathit{alice}, \mathit{carol})$; however, it
cannot infer $\texttt{grandparent}(\mathit{alice}, X)$ for a fresh variable
$X$ without performing unification.  The test
\texttt{prolog\_grandparent\_does\_not\_derive\_reversed} confirms that the
reversed query $\texttt{grandparent}(\mathit{carol}, \mathit{alice})$ is
\emph{not} derivable---a lookup table that stores both orderings would pass
both tests accidentally, but the
\texttt{prolog\_grandparent\_does\_not\_confuse\_parent\_with\_grandparent}
test introduces a sibling-step query that a table would conflate.  Together
the three tests require genuine SLD resolution: the shared-variable
mechanism is not expressible as a finite enumeration of ground facts.
\qquad $\square_{(ii)}$

\medskip\noindent
\textbf{(iii) $\STRIPS$ Plan Soundness \cite{fikes1971strips}.}
A $\STRIPS$ plan is a sequence of operators $\langle o_1, \ldots, o_k
\rangle$ such that the preconditions of each $o_j$ are satisfied by the
state obtained after applying $o_1, \ldots, o_{j-1}$ to the initial state.
The plan-soundness tests verify that the returned plan, when executed in the
simulator, reaches the goal state and that no intermediate state violates
a precondition.  A lookup table mapping $(s_0, g)$ to a fixed sequence
fails on any $(s_0', g')$ pair not present in the table; the planner's
domain requires $O(2^n)$ state-pairs for $n$ propositions, ruling out
any polynomial-size table.  \qquad $\square_{(iii)}$

\medskip\noindent
\textbf{(iv) $\SOAR$ Antisymmetry \cite{laird1987soar}.}
The $\SOAR$ breed encodes preference semantics: if operator $o_1$ is
preferred to $o_2$, then $o_2$ must \emph{not} be preferred to $o_1$.
The antisymmetry tests at lines 479--534 of
\texttt{strips\_soar\_cbr\_invariants.rs} assert $\neg\,\mathrm{pref}(o_2,
o_1)$ for every pair $(o_1, o_2)$ where $\mathrm{pref}(o_1, o_2)$ holds.
A lookup table that stores preference pairs without enforcing the
antisymmetry constraint will fail on the complement queries; enforcing
antisymmetry requires relational closure, which cannot be expressed as a
finite stored relation without re-deriving closures.  \qquad $\square_{(iv)}$

\medskip\noindent
\textbf{(v) CBR Jaccard Identity \cite{aamodt1994cbr}.}
The $\CBR$ breed computes case-based similarity using the Jaccard index
$J(A,B) = |A \cap B| / |A \cup B|$.  The identity tests at lines 755--767
of \texttt{strips\_soar\_cbr\_invariants.rs} assert $J(A, A) = 1$ for
all $A$ in the test corpus.  Any lookup table that hardcodes $J(A, A)$
only for the training cases will fail on a fresh case $A'$ not in the
table; moreover, a table entry for $J(A, A)$ cannot be reused to answer
$J(A, B)$ for $B \neq A$ without computing set intersection and union.
\qquad $\square_{(v)}$

\medskip\noindent
\textbf{(vi) $\HEARSAY$ KSAR Opportunistic Scheduling \cite{erman1980hearsay}.}
The $\HEARSAY$ breed implements a blackboard architecture in which
Knowledge Source Activation Records (KSARs) are scheduled opportunistically
by a scheduler that ranks them on a priority score.  The KSAR tests verify
that higher-priority activations are always selected before lower-priority
ones regardless of arrival order.  A lookup table over a fixed set of
activation sequences fails on any permutation of arrivals not seen during
construction; the scheduler's correctness is a property of the total order
on priority scores, which is algebraically defined and not enumerable.
\qquad $\square_{(vi)}$

\medskip\noindent
\textbf{(vii) $\DENDRAL$ Forbid Absoluteness \cite{feigenbaum1971dendral}.}
The $\DENDRAL$ breed applies a spectral interpretation filter that must
reject hypotheses whose predicted spectrum contains a peak forbidden by
the evidence.  The absoluteness falsification test at lines 348--376 of
\texttt{gps\_dendral\_eliza\_falsification.rs} asserts that a hypothesis
containing an absolutely forbidden substructure is rejected even when it
also satisfies many positive constraints.  A lookup table mapping
substructures to accept/reject fails on novel substructures; the forbidden
list is extensible and the test introduces a substructure outside the
training corpus.  \qquad $\square_{(vii)}$

\medskip\noindent
\textbf{(viii) $\GPS$ Difference-Table Coverage \cite{newell1963gps}.}
The $\GPS$ breed reduces goals by matching differences between current
state and goal state against a difference table that specifies which
operators reduce each type of difference.  The coverage tests verify that
every difference type in the test domain maps to a non-empty set of
applicable operators, and that the planner selects an operator from the
correct set.  A lookup table over difference--operator pairs fails when a
novel difference type is introduced; the GPS specification requires full
coverage of the difference-table schema, not just the training instances.
\qquad $\square_{(viii)}$

\medskip\noindent
\textbf{(ix) $\ELIZA$ Keystack Priority \cite{weizenbaum1966eliza}.}
The $\ELIZA$ breed implements a pattern-matching dialogue system in which
patterns are assigned priority scores and matched in priority order against
the input.  The keystack-priority tests in
\texttt{gps\_dendral\_eliza\_falsification.rs} assert that a higher-priority
pattern always preempts a lower-priority one, even when the lower-priority
pattern also matches the input.  A lookup table over fixed input strings
fails on inputs that match multiple patterns in different priority orders
than those seen in training; determining the correct response requires
evaluating all applicable patterns and selecting by priority, which is
procedural and not tabulatable.  \qquad $\square_{(ix)}$

\medskip\noindent
Clauses (i)--(ix) cover all thirteen breeds (the remaining four share the
algebraic structure of $\SOAR$ antisymmetry or $\PROLOG$ variable-binding
and are subsumed by those arguments).  Therefore the 483-test suite is an
unfakeable oracle.
\end{proof}

% ─────────────────────────────────────────────────────────────────────────────
\section{The Falsifiability Metric}
\label{sec:falsifiability-metric}
% ─────────────────────────────────────────────────────────────────────────────

Theorem~\ref{thm:oracle} establishes unfakeability qualitatively.  We now
introduce a quantitative measure of how \emph{much} evidence each test
suite provides.

\begin{definition}[Falsifiability Pseudometric]
\label{def:falsifiability-metric}
Let $\mathcal{T}_i$ be the test suite for breed $b_i$ and let
$\mathcal{A}_i$ denote the set of all algorithmically simpler candidate
functions (lookup tables, finite automata, polynomial interpolants, etc.).
Define the \emph{falsifiability pseudometric}
\[
  d_F(\mathcal{T}_i)
  \;=\;
  1 - \frac{\#\bigl\{\hat{f} \in \mathcal{A}_i \mid \hat{f} \text{ passes }
    \mathcal{T}_i\bigr\}}%
  {|\mathcal{A}_i|}.
\]
$d_F(\mathcal{T}_i) = 1$ means no simplified impostor survives; $d_F =
0$ means the suite exerts no selection pressure.  Because $\mathcal{A}_i$
is countably infinite we interpret the ratio as the limit of the density
over finite prefixes of a canonical enumeration of $\mathcal{A}_i$.
\end{definition}

The pseudometric $d_F$ satisfies $d_F \in [0,1]$, symmetry in the
trivial sense (it is a property of a single suite), and the triangle
inequality vacuously, hence the name ``pseudometric'' rather than
``metric''.  Its primary use is comparative: a suite $\mathcal{T}$ with
higher $d_F$ provides stronger falsifiability evidence.

\begin{corollary}[Operational Falsifiability of All Thirteen Breeds]
\label{cor:operational-falsifiability}
For each breed $b_i$ $(i = 1, \ldots, 13)$ the sub-suite $\mathcal{T}_i$
satisfies $d_F(\mathcal{T}_i) = 1$.
\end{corollary}

\begin{proof}
By Theorem~\ref{thm:oracle}, every simplified impostor $\hat{f} \in
\mathcal{A}_i$ fails at least one test in $\mathcal{T}_i$.  The numerator
in Definition~\ref{def:falsifiability-metric} is therefore zero, giving
$d_F = 1 - 0 = 1$.
\end{proof}

Corollary~\ref{cor:operational-falsifiability} is the Popperian content of
the framework: the 483-test suite maximally falsifies the hypothesis that
any of the thirteen breeds can be replaced by a simpler computational
artefact.  This is the operational meaning of the epigraph: scientific
status is conferred not by confirmation but by survival of the strongest
refutation attempt.

% ─────────────────────────────────────────────────────────────────────────────
\section{Summary}
\label{sec:oracle-summary}
% ─────────────────────────────────────────────────────────────────────────────

This chapter formalised the notion of a complete test suite
(Definitions~\ref{def:test}--\ref{def:unfakeable}) and proved that the
483-test \textsc{wasm4pm} suite constitutes an unfakeable oracle for all
thirteen breeds (Theorem~\ref{thm:oracle}).  The nine canonical witness
tests---spanning $\MYCIN$ boundary arithmetic, $\PROLOG$ shared-variable
unification, $\STRIPS$ plan soundness, $\SOAR$ antisymmetry, $\CBR$ Jaccard
identity, $\HEARSAY$ scheduling priority, $\DENDRAL$ forbidden-substructure
rejection, $\GPS$ difference-table coverage, and $\ELIZA$ keystack
preemption---collectively rule out every algorithmically simpler impostor.
The falsifiability pseudometric $d_F$ (Definition~\ref{def:falsifiability-metric})
achieves its maximum value of 1 for each breed (Corollary~\ref{cor:operational-falsifiability}).
Chapter~\ref{chap:ocel} now asks a complementary question: given that each
breed \emph{passes} its oracle, can we verify at runtime that every
\emph{execution} also conformed to the declared process model?  The answer
requires an object-centric event-log framework, which we develop next.
